\subsection{Primer alignment with the Smith-Waterman algorithm}
\label{sec:conclud}

\paragraph{Presentation}ConCluD is a software that relies on De-Buijn graph partitioning to generate consensus sequences in the context of DNA storage decoding~\cite{DeMoor2025}.
The partitioning step requires to align reads against a set of known primers to find the best fitting one.
The published version of ConCluD uses a simple approach based on the Hamming distance.
Such computation has already been investigated in Section~\ref{sec:mapping} with the MiMyCS software.
To explore a different scenario, we thus decide to rather re-implement with the BPL library an older version of ConCluD where the alignment step was originally done with the Smith-Waterman algorithm.

\paragraph{Limitations encountered} None: the code can be fully architecture-agnostic.

\paragraph{PiM specifics} To perform the alignment computations as efficiently as possible on the UPMEM architecture, we write manually a few lines of assembly code.
In particular, we take advantage of the \texttt{cmp4} instruction to compare four bytes at once.
This part of the task thus uses target-specific code.

\paragraph{Benchmarks} We use three real datasets from the experimental evaluation in ~\cite{DeMoor2025}, namely Odeuropa, Dunhuang-d0 and HorrorMovie.
For PiM, we use 32 ranks.

\paragraph{Results} We give results in \tablename~\ref{tab:sw-results}.
On all datasets, the BPL version is slower than the SDK version regarding the PiM architecture.
For the largest dataset, HorrorMovie, it is 34\% slower.
On top of the BPL overhead, another reason contributes to this observation: we had to split the work statically between tasklets to keep the program entirely architecture-agnostic with regard to the current BPL features.
On the other hand, the SDK version uses dynamic scheduling to assign reads to tasklets on the same DPU.
The latter is more efficient because the processing time of a read can differ drastically: some reads lead to a perfect alignment very fast, while other require to go through the whole Smith-Waterman matrix or even try several primers.

Regarding the Multicore target, the BPL version is also slower, except for the Dunhuang-d0 dataset where it is faster.
The fact that the BPL version processes reads in batches instead of one at a time, combined with the dataset having only one primer, seem to make the BPL version less memory-bound that the Standard program in this case.

\begin{table}[htbp]
    \centering
    \caption{Results of the primer alignment benchmark}
    \label{tab:sw-results}
    \begin{tabular}{lrrrr}
    \toprule
    & \multicolumn{2}{c}{Multicore time (s)} & \multicolumn{2}{c}{PiM time (s)} \\
    \cmidrule(l{3pt}r{3pt}){2-3}\cmidrule(l{3pt}r{3pt}){4-5}
    & Standard & BPL & SDK & BPL \\
    \midrule
    Odeuropa & 12.5 & 16.3 & 4.7 & 6.7 \\
    Dunhuang-d0 & 37.0 & 25.8 & 12.4 & 20.6 \\
    HorrorMovie & 262.4 & 382.0 & 92.3 & 123.9 \\
    \bottomrule
    \end{tabular}
\end{table}

\paragraph{Source code} The BPL implementation is publicly available at \url{https://gitlab.inria.fr/pim/org.pim.conclud-bpl} (GNU AGPLv3).
The original Standard Multicore and SDK implementations are available at \url{https://gitlab.inria.fr/pim/org.pim.dnarxiv/-/tags/experiments-bioRxiv-preprint} (GNU AGPLv3).